% Large Language Model (LLM) agents have demonstrated remarkable potentials in automating complex, open-ended tasks, such as software engineering~\citep{}, web navigation~\citep{} and deep research~\citep{}.
% An essential characteristics of these tasks is their long-horizon nature--as the number of interactions grows, the conversation history also scales linearly. This expansion not only risks exceeding the context window but also presses the computation overhead incurred due to long-context processing.
% Efficiently managing this context is therefore critical to satisfy the limits of hardware constraints and desired inference latency.
The capabilities of Large Language Model (LLM) agents have expanded significantly, moving beyond simple conversational assistants to autonomous systems capable of tackling complex, multi-step problems in domains such as software engineering~\citep{jimenez2023swe}, scientific discovery~\cite{lu2024aiscientistfullyautomated}, and open-ended exploration~\cite{he2024webvoyagerbuildingendtoendweb,tongyidr}. Unlike standard question-answering tasks, the success on these applications depends not only on the immediate instruction but on the agent's ability to maintain a coherent understanding of its entire interaction history. For instance, in repository-level code generation, an agent must recall decisions made thousands of steps prior to generate valid subsequent actions. Consequently, the ability to process and reason over extremely long contexts has become a non-negotiable requirement for high-performing agentic systems~\cite{zhou2025mem1,wu2025resum,lu2025scaling,sun2025scaling,yu2025memagent}.

% The cost of long-context processing typically can be split into two distinct phases: prefilling and decoding. While prefilling latency scales quadratically with the context length, decoding time usually dominates the end-to-end latency given that capable agent models usually generate extensive reasoning before proposing the policies. 
% Thus the most important 
Despite the increasing attention to extend the effective context window, there has been surprisingly little attention paid to the severe computational costs introduced by long-context processing during agentic inference. While modern LLMs can natively accept context windows spanning millions of tokens, utilizing them during online agent deployment is often prohibitively expensive~\citep{yuan2024llm}. The primary challenge lies in the inference latency of the decoding stage. As the interaction history grows, the KV cache footprint expands linearly, and the attention computation cost scales, saturating the High Bandwidth Memory (HBM) of modern GPUs. Unlike the prefilling phase, which is compute-bound and parallelizable, the auto-regressive decoding phase is memory-bound, since fetching the entire KV cache for every generated token results in high latency that degrades the user experience and limits the throughput of real-time systems~\citep{tang2024questqueryawaresparsityefficient}.
% This creates a "memory wall" where the agent's reasoning power is throttled by the sheer volume of its past observations~\citep{tang2024questqueryawaresparsityefficient}.

% One increasingly popular approach to alleviate this issue is to condense the interaction history into a compact summary for the agent to process~\citep{}.
% For instance, \cite{} summarizes 
% However, na\"ive summarization further burdens the 
To address these computational challenges, previous solutions have largely pursued two directions: context reduction and system-level optimization.
Context reduction strategies attempt to limit the number of tokens the model must attend to. Sliding-window attention~\citep{fu2025slidingwindowattentiontraining} restricts the agent to a fixed local context, inherently sacrificing the ability to recall distant dependencies, such as the initial user goal or a bug identified thousands of steps prior. Retrieval-Augmented Generation (RAG)~\cite{wu2025longmemevalbenchmarkingchatassistants,shi2026lookreasonforwardrevisitable} alleviates this by fetching relevant snippets from the history. However, RAG relies heavily on semantic similarity, which often fails to capture the high-level state or procedural trajectory of an agent, leading to fragmented reasoning.

On the other hand, system-level optimizations aim to make processing the full context more efficient without reducing the number of tokens. Modern inference engines like vLLM~\cite{kwon2023efficientmemorymanagementlarge} have revolutionized memory management through PagedAttention, significantly reducing fragmentation. Similarly, Sparse Attention mechanisms~\cite{xiao2024efficientstreaminglanguagemodels} and KV Cache Quantization~\cite{hooper2025kvquant10millioncontext} reduce the compute and memory bandwidth requirements by sparsifying interactions or lowering numerical precision.
% KV cache can also be offloaded to CPU for future decoding~\cite{jin2024compute}.
While these methods provide substantial throughput improvements, they do not solve the the redundancy of re-encoding or re-attending to a massive, largely static interaction history for every single decoding step. As a result, even highly optimized engines eventually hit a latency wall as the interaction trace grows indefinitely.

To bridge this gap, we introduce \model, a framework that decouples the distinct responsibilities of memory management and agentic reasoning. Unlike prior approaches that force a single large model to handle both history compression and policy generation, \model offloads the heavy lifting of long-context processing to a dedicated, lightweight summarization model. This architectural separation allows us to propose a novel $k$-step-off asynchronous pipeline, where the memory model continuously compresses history in the background, freeing the main agent to decode with a significantly reduced context window. To ensure this compressed state effectively guides the agent, we introduce a reward-driven alignment strategy that trains the memory model to capture the "sufficient statistics" required for optimal decision-making. By shifting the computational burden from the critical path of decoding to a parallelizable background process, \model achieves the best of both worlds: the global reasoning capability of long-context models and the low-latency inference of short-context systems.

Our contributions are summarized as follows:
\begin{itemize}
    \item We propose \model, a framework that separates memory management from reasoning, enabling the use of specialized, lightweight models for efficient history compression.
    \item Asynchronous Inference: We design a $k$-step-off pipeline that overlaps memory summarization with agent execution, masking the latency of context processing and reducing the decoding overhead by up to $1.4\times$.
    \item We introduce a novel training methodology that aligns the memory model using a functional equivalence reward, ensuring the compressed summary recovers the full-context agent's policy without expensive online interaction.
    \item Extensive experiments on SWE-Bench-Verified demonstrate that \model significantly reduces latency while maintaining competitive performance against state-of-the-art long-context baselines. Code is available at: \url{https://github.com/horizon-llm/CoMem.git}.
\end{itemize}